\documentclass[11pt]{article}
\usepackage{mathblatt}
\begin{document}
% Kopfzeile Seite 1 ohne Kennung; die Rückseite (Lösungen, für den Lehrer) bekommt sie nach \begleitteil
\blattkopf{Mathematik}{Basisaufgaben}
\einheitenkopf[][Zettel 19]{Basisaufgaben · Zettel 19}
\zweigzeile{je Aufgabe 1 Punkt · Lösungen auf der Rückseite}
\par\vspace{2pt}\noindent Name \underline{\hspace{7cm}}\hfill Datum \underline{\hspace{3cm}}\hfill Punkte \underline{\hspace{1.2cm}} / 18\par\vspace{6pt}
% \small: volle Seite, kurze Aufgaben zweispaltig (v0.6)
\small

\noindent\begin{minipage}[t]{0.485\linewidth}\raggedright
% 1: Bruchteil einer Größe berechnen – brueche-dezimalzahlen-basis-k2-v8 (Original 2018-OS-B1a, ausgedacht: keine Marke)
\begin{aufgabe}{$\frac{4}{5}$ von $2{,}5$ l – wie viel Liter? \leerfeld[l]}
\end{aufgabe}
\end{minipage}\hfill\begin{minipage}[t]{0.485\linewidth}\raggedright
% 2: Exponent einer Potenz bestimmen – potenzen-wurzeln-basis-k1-v7 (Original 2024-OS-B1g, ausgedacht: keine Marke)
\begin{aufgabe}{$6^x = 36$ – welche Zahl ist $x$? x = \leerfeld}
\end{aufgabe}
\end{minipage}\par

\noindent\begin{minipage}[t]{0.485\linewidth}\raggedright
% 3: Grundwert berechnen – prozentrechnung-basis-k1-v6 (Original 2025-OS-B1a, ausgedacht: keine Marke)
\begin{aufgabe}{Beim Kauf einer Lampe spart Ole $15\,€$, das sind $75\,\%$ des alten Preises. Wie hoch war der alte Preis? \leerfeld[€]}
\end{aufgabe}
\end{minipage}\hfill\begin{minipage}[t]{0.485\linewidth}\raggedright
% 4: Zeiteinheiten umrechnen – einheiten-basis-k5-v5 (Original 2025-OS-B1b, ausgedacht: keine Marke)
\begin{aufgabe}{Eine Ampel bleibt $3$ min rot. Gib die Zeit in Sekunden an. \leerfeld[s]}
\end{aufgabe}
\end{minipage}\par

\noindent\begin{minipage}[t]{0.485\linewidth}\raggedright
% 5: Term durch Zusammenfassen gleichartiger Glieder vereinfachen – terme-basis-k1-v3 (Original 2025-GYM-B2a, ausgedacht: keine Marke)
\begin{aufgabe}{Vereinfache den Term $5a + 2a^2 - 8a$ und berechne seinen Wert für $a = 2$.}
\end{aufgabe}
\end{minipage}\hfill\begin{minipage}[t]{0.485\linewidth}\raggedright
% 6: Lineare Gleichung lösen – lineare-gleichungen-basis-k1-v12 (Original 2024-OS-B1d, ausgedacht: keine Marke)
\begin{aufgabe}{Löse die Gleichung $3 \cdot (x + 2) = 21$. x = \leerfeld}
\end{aufgabe}
\end{minipage}\par

\noindent\begin{minipage}[t]{0.485\linewidth}\raggedright
% 7: Rechteckseite aus Umfang berechnen – flaechen-basis-k1-v3 (Original 2018-OS-B1f, ausgedacht: keine Marke)
\begin{aufgabe}{Ein rechteckiger Spielplatz ist mit $40$ m Zaun eingefasst und $12$ m lang. Wie breit ist er? \leerfeld[m]}
\end{aufgabe}
\end{minipage}\hfill\begin{minipage}[t]{0.485\linewidth}\raggedright
% 8: Trigonometrische Gleichung nach Seite umstellen – trigonometrie-basis-k1-v3 (Original 2020-OS-B1j, ausgedacht: keine Marke)
\begin{aufgabe}{$\mathrm{tan}\,45^\circ = \frac{11}{x}$ – berechne $x$. x = \leerfeld}
\end{aufgabe}
\end{minipage}\par

\noindent\begin{minipage}[t]{0.485\linewidth}\raggedright
% 9: Arithmetisches Mittel berechnen – daten-basis-k1-v6 (Original 2021-OS-B1f, ausgedacht: keine Marke)
\begin{aufgabe}{Drei Hunde wiegen $12$; $30$; $18$ kg. Berechne das arithmetische Mittel (Masse). \leerfeld[kg]}
\end{aufgabe}
\end{minipage}\hfill\begin{minipage}[t]{0.485\linewidth}\raggedright
% 10: Median bestimmen – daten-basis-k3-v4 (Original 2024-OS-B1h, ausgedacht: keine Marke)
\begin{aufgabe}{Preise für eine Kugel Eis in Euro: $1{,}50$; $1{,}20$; $1{,}80$; $1{,}40$. Bestimme den Median. \leerfeld[€]}
\end{aufgabe}
\end{minipage}\par

% 11: Spannweite berechnen – daten-basis-k4-v2 (Original 2019-OS-B1i, ausgedacht: keine Marke)
\begin{aufgabe}{Fünf Kinder sind so groß: $158$; $171$; $164$; $149$; $166$ cm. Gib die Spannweite an. \leerfeld[cm]}
\end{aufgabe}

% 12: Strecke aus Teilstrecken berechnen – einheiten-basis-k4-v2 (Original 2018-OS-B1b, ausgedacht: keine Marke)
\begin{aufgabe}{Eine Schnur ist $2{,}5$ m lang. Davon werden $6$ Stücke von je $25$ cm abgeschnitten. Wie lang ist der Rest? \leerfeld[cm]}
\end{aufgabe}

% 13: Scheitelpunktform aufstellen – quadratische-funktionen-basis-k4-v2 (Original 2016-OS-B1g, ausgedacht: keine Marke)
\begin{aufgabe}{Eine verschobene Normalparabel hat den Scheitelpunkt $S(-4|3)$. Welche Gleichung gehört zu ihr? Kreuze an.\\ \kreuz{$y = (x + 4)^2 - 3$}\\ \kreuz{$y = (x - 3)^2 - 4$}\\ \kreuz{$y = (x + 4)^2 + 3$}\\ \kreuz{$y = (x - 4)^2 + 3$}}
\end{aufgabe}

% 14: Wahrscheinlichkeit einstufig – bruchrechnung-basis-k1-v8 (Original 2016-OS-B1f, ausgedacht: keine Marke)
\begin{aufgabe}{In einer Tüte sind 6 rote und 2 gelbe Bonbons. Lea nimmt ohne hinzusehen eines. Wie groß ist die Wahrscheinlichkeit für ein gelbes? \leerfeld}
\end{aufgabe}

% 15: Zufallsgerät zu Wahrscheinlichkeit entwerfen – wahrscheinlichkeit-basis-k2-v5 (Original 2019-OS-B1g, ausgedacht: keine Marke)
\begin{aufgabe}{Ein Glücksrad hat $10$ gleich große Felder. Die Wahrscheinlichkeit für einen Gewinn soll $60\,\%$ sein. Wie viele Gewinnfelder braucht man? \leerfeld Felder}
\end{aufgabe}

% 16: Satz des Pythagoras formulieren – pythagoras-basis-k2-v2 (Original 2022-OS-B1g, ausgedacht: keine Marke)
\begin{aufgabe}{Im Dreieck $ABC$ liegt der rechte Winkel bei $A$. Welche Gleichung gilt? Kreuze an. \\ \kreuz{$b^2 + c^2 = a^2$} \kreuz{$a^2 + b^2 = c^2$} \kreuz{$b \cdot c = a^2$}}
\end{aufgabe}

% 17: Parabel verschieben – quadratische-funktionen-basis-k2-v1 (Original 2015-OS-B1i, ausgedacht: keine Marke)
\begin{aufgabe}{\begin{minipage}[t]{0.6\linewidth}Die Normalparabel wird um $3$ Einheiten nach links verschoben. Gib den Scheitelpunkt der neuen Parabel an. S( \leerfeld | \leerfeld )\end{minipage}\hfill\begin{minipage}[t]{0.37\linewidth}\vspace{0pt}
\begin{ksys}[xmin=-5,xmax=5,ymin=-5,ymax=5,klein] \parabel{1}{0}{0}{} \end{ksys}
\end{minipage}}
\end{aufgabe}

% 18: Winkel im Viereck berechnen – winkel-dreiecke-basis-k4-v3 (Original 2026-FOR-B1i, ausgedacht: keine Marke)
\begin{aufgabe}{\begin{minipage}[t]{0.6\linewidth}Ein Parallelogramm hat links unten einen Winkel von $72^\circ$. Wie groß ist der Winkel $\gamma$ rechts oben? $\gamma =$ \leerfeld °\end{minipage}\hfill\begin{minipage}[t]{0.37\linewidth}\vspace{0pt}
\parallelogramm[winkel={72^\circ,,\gamma,}]{3}{1.8}{72}
\end{minipage}}
\end{aufgabe}

\begleitteil
\blattkopf{Basisaufgaben · Zettel 19}{BAS-Z19 · Lösungen}
\erg{1}{$2{,}5 : 5 \cdot 4 = 2$ l \hfill {\footnotesize\textit{Bruchteil einer Größe berechnen · 2018}}}
\erg{2}{$6, 36$: zwei Faktoren, also $x = 2$ \hfill {\footnotesize\textit{Exponent einer Potenz bestimmen · 2024}}}
\erg{3}{$15 : 75 \cdot 100 = 20\,€$ (nicht $75\,\%$ von $15\,€$) \hfill {\footnotesize\textit{Grundwert berechnen · 2025}}}
\erg{4}{$3 \cdot 60 = 180$ s \hfill {\footnotesize\textit{Zeiteinheiten umrechnen · 2025}}}
\erg{5}{$5a - 8a = -3a$, Term $2a^2 - 3a$; für $a = 2$: $2 \cdot 2^2 - 3 \cdot 2 = 8 - 6 = 2$ \hfill {\footnotesize\textit{Term durch Zusammenfassen gleichartiger Glieder vereinfachen · 2025}}}
\erg{6}{$x + 2 = 7$, also $x = 5$ \hfill {\footnotesize\textit{Lineare Gleichung lösen · 2024}}}
\erg{7}{$40 : 2 = 20$; $20 - 12 = 8$ m \hfill {\footnotesize\textit{Rechteckseite aus Umfang berechnen · 2018}}}
\erg{8}{$x = 11 : \mathrm{tan}\,45^\circ = 11 : 1 = 11$ \hfill {\footnotesize\textit{Trigonometrische Gleichung nach Seite umstellen · 2020}}}
\erg{9}{$60 : 3 = 20$ kg \hfill {\footnotesize\textit{Arithmetisches Mittel berechnen · 2021}}}
\erg{10}{geordnet $1{,}20$; $1{,}40$; $1{,}50$; $1{,}80$; $(1{,}40 + 1{,}50) : 2 = 1{,}45$ € \hfill {\footnotesize\textit{Median bestimmen · 2024}}}
\erg{11}{größter minus kleinster Wert: $171 - 149 = 22$ cm \hfill {\footnotesize\textit{Spannweite berechnen · 2019}}}
\erg{12}{$2{,}5$ m $= 250$ cm; $250 - 6 \cdot 25 = 250 - 150 = 100$ cm \hfill {\footnotesize\textit{Strecke aus Teilstrecken berechnen · 2018}}}
\erg{13}{$y = (x + 4)^2 + 3$ \hfill {\footnotesize\textit{Scheitelpunktform aufstellen · 2016}}}
\erg{14}{$\frac{2}{8} = \frac{1}{4}$ \hfill {\footnotesize\textit{Wahrscheinlichkeit einstufig · 2016}}}
\erg{15}{$10 \cdot 0{,}6 = 6$ \hfill {\footnotesize\textit{Zufallsgerät zu Wahrscheinlichkeit entwerfen · 2019}}}
\erg{16}{$b^2 + c^2 = a^2$ \hfill {\footnotesize\textit{Satz des Pythagoras formulieren · 2022}}}
\erg{17}{$S(0|0)$ wandert nach $S(-3|0)$ \hfill {\footnotesize\textit{Parabel verschieben · 2015}}}
\erg{18}{$\gamma = 72^\circ$ (gegenüberliegende Winkel sind gleich groß) \hfill {\footnotesize\textit{Winkel im Viereck berechnen · 2026}}}
\end{document}
